Chương này đóng vai trò thiết lập nền tảng lý luận cho toàn bộ khóa luận, tổng quan về bối cảnh nghiên cứu và các cơ sở khoa học cốt lõi dẫn dắt đến hướng tiếp cận của đề tài. Đề tài xoay quanh sự hội tụ đa ngành: kết hợp chặt chẽ giữa kỹ thuật giám sát vi khí hậu hiện đại, công nghệ thủy canh sinh thái (với trọng tâm là cây lan ý) và các hệ thống tự động hóa. Tất cả các thành tố này được ứng dụng trong một mạng lưới Internet vạn vật (IoT) bảo đảm tính đồng bộ, độ trễ thấp và khả năng vận hành tự chủ ở quy mô cục bộ.

\section{Đặt vấn đề}

Trong kỷ nguyên đô thị hóa hiện đại, chất lượng không khí trong không gian kín (IAQ) đã vượt ra khỏi khuôn khổ của những tiện nghi , trở thành một yếu tố có tác động trực tiếp và sâu sắc đến sức khỏe cộng đồng, năng lực nhận thức cũng như trạng thái cân bằng sinh học của con người. Các báo cáo  uy tín liên tục cảnh báo rằng, sự suy giảm chất lượng không khí tại các không gian kín, kết hợp với chỉ số lưu thông khí tươi thấp, là nguyên nhân sâu xa kích hoạt và làm trầm trọng thêm các bệnh lý hô hấp mãn tính, hội chứng bệnh nhà kín (Sick Building Syndrome), và suy giảm chất lượng sống tổng thể \cite{who_air_pollution,epa_iaq}. 

Nhìn vào thực tế kiến trúc của các căn hộ chung cư đương đại, chúng ta dễ dàng nhận thấy xu hướng tối đa hóa diện tích sử dụng bằng sự khép kín . Tuy nhiên,sự khép kín này lại tạo ra sự phụ thuộc vào các hệ thống điều hòa không khí, quạt thông gió cưỡng bức . Khi mật độ sinh hoạt của con người gia tăng trong các không gian hạn hẹp, tốc độ tích tụ của khí CO$_2$ diễn ra cực kỳ nhanh chóng. Kéo theo đó là sự xáo trộn liên tục của  nhiệt độ và độ ẩm dưới tác động của thiết bị cơ điện. Điều đáng quan ngại là những biến thiên vi khí hậu này thường diễn tiến một cách âm thầm, vượt qua ngưỡng dung sai và độ nhạy cảm của các giác quan thông thường, khiến con người mất đi khả năng nhận diện và can thiệp kịp thời.

Đứng trước thách thức đó, việc đưa các mảng xanh thực vật vào không gian sống đang thu hút sự quan tâm lớn của giới học thuật, không chỉ dưới lăng kính thẩm mỹ kiến trúc mà còn như một cỗ máy lọc khí sinh thái đầy hứa hẹn. Trong số các loài thực vật nội thất, lan ý (\textit{Spathiphyllum}) đặc biệt nổi bật nhờ phổ thích ứng sinh lý cực rộng. Loài cây này có khả năng quang hợp hiệu quả ở cường độ sáng thấp, chịu đựng tốt các cú sốc nhiệt ẩm, và đặc biệt là sinh trưởng cực kỳ ổn định trong môi trường thủy canh \cite{peace_lily_care_bhg,peace_lily_hydro_brightlane}. Các nghiên cứu chuyên sâu về sinh lý thực vật đã làm sáng tỏ khả năng của lan ý trong việc đồng hóa CO$_2$ và chủ động hấp thụ các hợp chất hữu cơ dễ bay hơi (VOC) độc hại—như benzene, formaldehyde, hay trichloroethylene—vốn rất phổ biến trong vật liệu nội thất \cite{hydroponic_iaq_survey,nasa_clean_air}. 

Dù vậy, khi xét dưới góc độ của hệ thống, năng lực xử lý sinh học của thực vật không phải là một hằng số. Nó là hàm phụ thuộc vào các biến số vi khí hậu như: động thái trao đổi chất tại vùng rễ, cường độ bức xạ quang hợp hiệu dụng và trạng thái hóa lý của dung dịch dinh dưỡng. Nhận thức sâu sắc được thực tiễn này, nghiên cứu của em không nhìn nhận lan ý như một bộ lọc khí thụ động để thay thế hoàn toàn các thiết bị cơ điện. Thay vào đó, em định vị hệ sinh thái thực vật này như một thành phần chủ động và được tinh chỉnh liên tục bởi một hệ thống kiểm soát môi trường tự động hóa.

Xuất phát từ những luận điểm trên, khóa luận  đề xuất một hệ thống IoT đa mục tiêu: vừa theo dõi sát sao sự biến thiên của chất lượng không khí trong không gian sống, vừa giám sát chặt chẽ các thông số sinh lý của cây thủy canh. Về mặt kiến trúc vật lý, hệ thống là một mạng lưới các cảm biến thông minh sử dụng vi điều khiển Arduino Pro Mini, giao tiếp không dây qua giao thức tầm xa LoRa để truyền dữ liệu về gateway. Tại đây, luồng dữ liệu được định tuyến lên nền tảng đám mây ThingsBoard nhằm mục đích lưu trữ dài hạn, phân tích trực quan hóa và thiết lập cổng giao tiếp với ứng dụng di động Flutter dành cho người dùng cuối. Không dừng lại ở việc quan sát thụ động, giá trị cốt lõi của đề tài nằm ở việc tích hợp cơ chế phản hồi chủ động. Thông qua các rơ-le điều khiển quạt, đèn và module phát hồng ngoại điều khiển điều hòa nhiệt độ, em đã hiện thực hóa một vòng lặp điều khiển khép kín: đo lường chính xác, cảnh báo kịp thời, và tự động kích hoạt các kịch bản can thiệp thích nghi nhằm duy trì trạng thái vi khí hậu tối ưu nhất.